% =============================================================================
%                         WEEK 1: INTRODUCTION
% =============================================================================
\section{Week 1: Introduction}

\subsection{Topic Summary}

\begin{keybox}[Learning Objectives]
By the end of this week, you will be able to:
\begin{itemize}
  \item Describe the goals and challenges of the Computer Vision field.
  \item Recognise real-world scenarios where Computer Vision is used.
  \item Identify and list various Computer Vision applications.
\end{itemize}
\end{keybox}

\textbf{Big Picture.}
This week introduces the field of Computer Vision: what it is, why it matters, why it's hard, and how we approach solving vision problems. Vision is our primary sense for experiencing the world, yet teaching machines to ``see'' remains one of the great challenges in AI. The key insight is that vision is an \term{ill-posed inverse problem}---many possible 3D scenes can produce the same 2D image. Solving this requires \term{priors} (constraints/expectations) that make some interpretations more likely than others.

\separator

\textbf{What Examiners Like to Ask:}
\begin{itemize}
  \item Define computer vision and distinguish it from related fields (image processing, computer graphics).
  \item Explain why vision is an ``ill-posed inverse problem.''
  \item Give examples of priors and how they constrain interpretations.
  \item List and categorise CV applications by type of information extracted.
  \item Explain why deep learning is not always the best solution.
\end{itemize}

% =============================================================================
\subsection{Core Concepts}

\textbf{Roadmap:}
\begin{enumerate}
  \item What is Computer Vision? (definitions, related fields)
  \item Why study vision? (importance, applications)
  \item Why is vision difficult? (ill-posed, exponentially large)
  \item How do we tackle the problem? (priors, engineering approaches)
\end{enumerate}

\bigseparator

% -----------------------------------------------------------------------------
\subsubsection{What is Computer Vision?}

\begin{defbox}[Computer Vision]
Computer Vision is a field of study that focuses on creating methods capable of extracting information about the world from images and videos.
\end{defbox}

\textbf{Classic definitions:}
\begin{itemize}
  \item ``Computing properties of the 3D world from one or more digital images'' (Trucco \& Verri)
  \item ``Making useful decisions about real physical objects and scenes based on sensed images'' (Stockman \& Shapiro)
  \item ``Extracting descriptions of the world from pictures or sequences of pictures'' (Forsyth \& Ponce)
\end{itemize}

\textbf{Alternative names:} machine vision, image analysis, image understanding, computational vision.

\separator

\textbf{Related Disciplines and Their Transformations:}

\begin{center}
\renewcommand{\arraystretch}{1.3}
\begin{tabular}{@{} l c c c @{}}
\toprule
\textbf{Field} & \textbf{Input} & $\rightarrow$ & \textbf{Output} \\
\midrule
Image Processing & Image & $\rightarrow$ & Image \\
Computer Vision & Image & $\rightarrow$ & Description \\
Computer Graphics & Description & $\rightarrow$ & Image \\
\bottomrule
\end{tabular}
\end{center}

\begin{tipbox}[Exam Tip: Field Distinctions]
A common question asks you to label arrows in a diagram showing transitions between ``Real World,'' ``Image,'' and ``Description.'' Remember:
\begin{itemize}
  \item Real World $\rightarrow$ Image: \textbf{Imaging/Photography}
  \item Image $\rightarrow$ Description: \textbf{Computer Vision}
  \item Description $\rightarrow$ Image: \textbf{Computer Graphics}
  \item Image $\rightarrow$ Image: \textbf{Image Processing}
\end{itemize}
\end{tipbox}

\textbf{Other related fields:}
\begin{itemize}
  \item \textbf{Pattern Recognition}: recognising and classifying stimuli in images and other datasets.
  \item \textbf{Photogrammetry}: obtaining measurements from images.
  \item \textbf{Biological Vision}: understanding visual perception in humans/animals (neuroscience, psychology, psychophysics).
\end{itemize}

\bigseparator

% -----------------------------------------------------------------------------
\subsubsection{Why Study Computer Vision?}

\textbf{Motivation:}
\begin{itemize}
  \item Vision is the main way we experience the world.
  \item Want machines to interact with and understand the world (robotics, AI).
  \item Digital images are everywhere---lots of information to extract.
  \item Many practical applications with real-world impact.
\end{itemize}

\separator

\textbf{Applications by Information Extracted:}

\begin{center}
\renewcommand{\arraystretch}{1.3}
\begin{tabular}{@{} p{4.5cm} p{8cm} @{}}
\toprule
\textbf{Application} & \textbf{Information Extracted} \\
\midrule
OCR / Character Recognition & Category of each letter/digit \\
Face Tracking (webcam) & Location of face in image \\
Face Recognition / Biometrics & Identity of individual \\
Content-Based Image Retrieval & Category of items, relationships between items \\
Driver Assistance & Categorisation (pedestrians, road edge) + 3D location \\
3D Modelling & 3D locations of all image points \\
Object Detection & Location + class of objects \\
Object Segmentation & Pixel-level object boundaries \\
Action Recognition & Temporal classification of actions \\
Visual SLAM & Camera pose + 3D map of environment \\
\bottomrule
\end{tabular}
\end{center}

\separator

\textbf{Application Categories:}
\begin{itemize}
  \item \textbf{Recognition}: OCR, face recognition, iris/fingerprint biometrics, landmark recognition
  \item \textbf{Detection \& Tracking}: face detection, people tracking, object tracking (e.g., Hawk-Eye in sports)
  \item \textbf{Retrieval}: content-based image retrieval, reverse image search
  \item \textbf{3D Reconstruction}: 3D models from images, structure from motion
  \item \textbf{Autonomous Systems}: driver assistance, self-driving cars, space exploration (Mars rovers)
  \item \textbf{Medical}: analysis of MRI, CT, ultrasound; surgical robotics
  \item \textbf{Entertainment}: AR/VR, TikTok filters, eye contact rectification
  \item \textbf{Industrial}: quality control inspection, barcode/QR scanning
\end{itemize}

\bigseparator

% -----------------------------------------------------------------------------
\subsubsection{Why is Computer Vision Difficult?}

\begin{thmbox}[Two Major Challenges]
\begin{enumerate}
  \item \textbf{One image $\rightarrow$ many interpretations}: the problem is \term{ill-posed}.
  \item \textbf{One object $\rightarrow$ many images}: the problem is \term{exponentially large}.
\end{enumerate}
\end{thmbox}

\separator

\textbf{Challenge 1: Vision is an Ill-Posed Inverse Problem}

\begin{defbox}[Forward vs Inverse Problems]
\begin{itemize}
  \item \textbf{Forward problem}: know causes, predict outcomes (well-posed, unique solution).
  \item \textbf{Inverse problem}: know outcomes, infer causes (often ill-posed, multiple solutions).
\end{itemize}
\end{defbox}

\begin{itemize}
  \item \textbf{Imaging} (forward): 3D world $\rightarrow$ 2D image is \textbf{unique} (well-posed).
  \item \textbf{Vision} (inverse): 2D image $\rightarrow$ 3D world is \textbf{NOT unique} (ill-posed).
\end{itemize}

For any given image, many different 3D scenes could have generated it. Humans typically perceive only one interpretation because our brains apply \term{priors} that make some interpretations more likely.

\begin{exbox}[Classic Ambiguous Figures]
\begin{itemize}
  \item \textbf{Necker Cube}: two equally valid 3D interpretations; perception flips spontaneously.
  \item \textbf{Rubin's Vase}: face or vase? Two interpretations, never perceived simultaneously.
  \item \textbf{Convex/Concave Craters}: prior expectation of light from above determines interpretation.
\end{itemize}
\end{exbox}

\separator

\textbf{Challenge 2: Vision Scales Exponentially}

Consider recognising an object that can appear at:
\begin{itemize}
  \item $l$ locations in the image
  \item $s$ different scales (sizes)
  \item $o$ orientations
  \item $c$ colours
\end{itemize}

This single object can generate $l \times s \times o \times c$ different images. The number of possible images grows \textbf{exponentially} with the number of parameters.

\separator

\textbf{Factors Affecting Appearance:}

\begin{center}
\renewcommand{\arraystretch}{1.3}
\begin{tabular}{@{} l p{9cm} @{}}
\toprule
\textbf{Factor} & \textbf{Effect} \\
\midrule
Viewpoint & Same object from different angles looks very different \\
Illumination & Different lighting changes appearance drastically \\
Non-rigid deformation & Object can bend, stretch (e.g., faces, animals) \\
Within-category variation & Objects in same category look different (e.g., different chairs) \\
Background clutter & Multiple objects in scene \\
Occlusion & Objects partially hidden by other objects \\
\bottomrule
\end{tabular}
\end{center}

\begin{tipbox}[Exam Tip: Discrimination vs Recognition]
``Objects that look very similar can be represented and recognized as different objects, whereas objects that look very different can be recognized as the same basic-level objects'' (Bar, 2004). This captures the dual challenge: distinguishing similar things while grouping dissimilar things.
\end{tipbox}

\separator

\textbf{Why Current Systems Are Limited:}
\begin{itemize}
  \item Most CV systems work in \textbf{specific domains}: specific task + specific environment.
  \item Solutions are \textbf{not robust} and \textbf{do not generalise} well.
  \item Code that works for one task may fail for similar tasks or same task under different conditions.
  \item The challenge: develop \textbf{general-purpose vision systems} that match human performance.
\end{itemize}

\begin{exbox}[Challenging Examples]
\begin{itemize}
  \item \textbf{Chihuahua or blueberry muffin?} Texture similarity fools classifiers.
  \item \textbf{Sheepdog or mop?} Shape/texture confusion.
  \item \textbf{How many cows?} Segmentation fails with unusual poses/occlusions.
\end{itemize}
These examples show that even state-of-the-art deep learning struggles with adversarial or unusual cases.
\end{exbox}

\bigseparator

% -----------------------------------------------------------------------------
\subsubsection{How Do We Tackle the Problem of Vision?}

\begin{defbox}[Prior (Constraint)]
A \term{prior} is additional information that helps reduce the number of possible solutions, making some interpretations more likely than others. Priors are essential for solving the ill-posed inverse problem of vision.
\end{defbox}

\textbf{Sources of Priors:}
\begin{enumerate}
  \item \textbf{Prior knowledge / familiarity}:
  \begin{itemize}
    \item Learned familiarity with certain objects
    \item Knowledge of image formation process (e.g., perspective, shadows)
  \end{itemize}
  \item \textbf{Prior exposure / motion / priming}:
  \begin{itemize}
    \item Recent or preceding sensory input
    \item Change blindness: disrupting temporal continuity hides changes
  \end{itemize}
  \item \textbf{Current context}:
  \begin{itemize}
    \item Surrounding visual scene
    \item Concurrent input from other senses
  \end{itemize}
\end{enumerate}

\separator

\textbf{Effects of Priors on Human Perception:}

\begin{center}
\renewcommand{\arraystretch}{1.3}
\begin{tabular}{@{} l p{9cm} @{}}
\toprule
\textbf{Prior Type} & \textbf{Example Effect} \\
\midrule
Illumination & Checker shadow illusion: squares A and B appear different but are same intensity (shadow prior) \\
Light direction & Craters appear convex/concave depending on assumed light source (from above) \\
Perspective & Ames room: people appear to change size due to perspective expectations \\
Prior knowledge & Dalmatian dog: invisible until you know what to look for \\
Context & ``THE CAT''---middle letter is H in first word, A in second (same shape) \\
Face detection & We see faces everywhere (Virgin Mary toast sold for \$28,000!) \\
\bottomrule
\end{tabular}
\end{center}

\begin{thmbox}[Illusions Reveal Assumptions]
Illusions are not ``mistakes'' by the visual system. They reveal the \textbf{assumptions/priors} the visual system uses to solve the under-constrained problem of vision. These assumptions work well in typical situations because they reflect how the world usually is.
\end{thmbox}

\separator

\textbf{Two Engineering Approaches:}

\begin{center}
\renewcommand{\arraystretch}{1.3}
\begin{tabular}{@{} l p{10cm} @{}}
\toprule
\textbf{Approach} & \textbf{Description} \\
\midrule
Forward Engineering & Determine requirements $\rightarrow$ design system $\rightarrow$ implement and refine. ``Top-down'': start with computational theory, fill in details. \\
Reverse Engineering & Find a system that works (e.g., brain) $\rightarrow$ analyse how it works $\rightarrow$ implement using same mechanisms. ``Bottom-up'': start with mechanisms, build model. \\
\bottomrule
\end{tabular}
\end{center}

This course considers \textbf{both} approaches:
\begin{itemize}
  \item \textbf{Computational (machine) vision}: How can we get computers to see?
  \item \textbf{Biological (human) vision}: How do people see?
\end{itemize}

\begin{tipbox}[Exam Tip: Deep Learning Is Not Always the Answer]
The lecture emphasises that traditional/classical CV methods sometimes outperform deep learning. Example: for endoscope content area segmentation, classical methods achieved better accuracy than DL approaches. Always consider the problem domain and data availability.
\end{tipbox}

\bigseparator

% =============================================================================
\subsection{Worked Examples}

\begin{exbox}[Example 1: Defining Computer Vision]
\textbf{Q:} Give a definition of ``Computer Vision.''

\textbf{A:} Computer Vision is a field of study that focuses on creating methods capable of extracting information about the world from images and videos.
\end{exbox}

\begin{exbox}[Example 2: Field Transitions Diagram]
\textbf{Q:} Label each arrow with the field that deals with that transition:
\begin{center}
Real World $\leftrightarrow$ Image $\leftrightarrow$ Description
\end{center}

\textbf{A:}
\begin{itemize}
  \item Real World $\rightarrow$ Image: \textbf{Imaging/Photography}
  \item Image $\rightarrow$ Real World: (Physical interaction, not a field)
  \item Image $\rightarrow$ Description: \textbf{Computer Vision}
  \item Description $\rightarrow$ Image: \textbf{Computer Graphics}
  \item Image $\rightarrow$ Image: \textbf{Image Processing}
\end{itemize}
\end{exbox}

\begin{exbox}[Example 3: Information Extracted by Application]
\textbf{Q:} What type of information is extracted in each application?

\textbf{A:}
\begin{itemize}
  \item \textbf{OCR}: category of each letter
  \item \textbf{Face tracking webcam}: location of face in image
  \item \textbf{Face recognition/biometrics}: identity of individual
  \item \textbf{Content-based image retrieval}: category of items, relationships between items
  \item \textbf{Driver assistance}: categorisation (pedestrians, road edge) + 3D location
  \item \textbf{3D modelling}: 3D locations of all image points
\end{itemize}
\end{exbox}

\begin{exbox}[Example 4: Ill-Posed Inverse Problem]
\textbf{Q:} Vision is often described as an ``ill-posed, inverse problem.'' Explain these terms and their opposites.

\textbf{A:}
\begin{itemize}
  \item \textbf{Forward problem}: Know the causes, predict the outcomes (e.g., know forces, calculate acceleration).
  \item \textbf{Inverse problem}: Know the outcomes, infer the causes (e.g., know pixel intensities, infer 3D scene).
  \item \textbf{Well-posed problem}: Has one unique solution.
  \item \textbf{Ill-posed problem}: Has multiple solutions (or no solution).
\end{itemize}
Vision is an \textbf{inverse problem} because we know pixel intensities (outcomes) and want to infer the 3D scene (causes). It is \textbf{ill-posed} because multiple 3D scenes can produce the same 2D image.
\end{exbox}

\begin{exbox}[Example 5: Priors in Vision]
\textbf{Q:} What is a ``prior'' and how does it help solve the ill-posed inverse problem of vision? Give examples.

\textbf{A:} A prior is additional information that reduces the number of possible solutions, making some interpretations more likely. Examples include:
\begin{itemize}
  \item \textbf{Knowledge of the world}: objects are usually solid, light comes from above
  \item \textbf{Image formation knowledge}: perspective geometry, how shadows work
  \item \textbf{Prior exposure}: recently seen objects are more easily recognised
  \item \textbf{Context}: surrounding scene constrains possible object identities
\end{itemize}
\end{exbox}

\bigseparator

% =============================================================================
\subsection{Summary}

\begin{keybox}[Key Takeaways]
\begin{enumerate}
  \item \textbf{Computer Vision} extracts descriptions/information from images (Image $\rightarrow$ Description).
  \item Vision is \textbf{difficult} because it is:
  \begin{itemize}
    \item \textbf{Ill-posed}: one image can have many 3D interpretations
    \item \textbf{Exponentially large}: one object can generate many images
  \end{itemize}
  \item \textbf{Priors/constraints} are essential for narrowing down interpretations.
  \item Human vision uses priors extensively---illusions reveal these assumptions.
  \item We tackle CV using both \textbf{forward engineering} (design systems) and \textbf{reverse engineering} (study biological vision).
  \item Deep learning is powerful but \textbf{not always the best solution}---classical methods can outperform in specific domains.
\end{enumerate}
\end{keybox}

\begin{center}
\renewcommand{\arraystretch}{1.3}
\begin{tabular}{@{} l l l @{}}
\toprule
\textbf{Concept} & \textbf{Definition} & \textbf{Implication} \\
\midrule
Inverse problem & Infer causes from outcomes & Multiple solutions possible \\
Ill-posed & No unique solution & Need priors to constrain \\
Prior & Additional constraining info & Makes some interpretations likely \\
Forward engineering & Top-down design & Computational theory first \\
Reverse engineering & Bottom-up modelling & Study biological systems \\
\bottomrule
\end{tabular}
\end{center}
